\section{Introduction}
\label{sec:intro}

The proliferation of task-specific deep neural networks and parameter-efficient fine-tuning (PEFT) adapters has spurred a major paradigm shift in multi-task learning. Instead of training massive, monolithic models from scratch on joint datasets, a highly successful and resource-efficient alternative is to train individual, specialized expert models and subsequently merge or ensemble them into a single unified network \cite{wortsman22, ilharco22, yadav23, yu23, shen24}. This practice, commonly referred to as model merging or model ensembling, eliminates catastrophic forgetting, maintains expert-level specialized performance, and avoids the computational and privacy hurdles of joint multi-task training.

Despite their empirical success, existing model merging and ensembling methods suffer from a severe, fundamental limitation: they operate under the implicit assumption of a \emph{flat Euclidean geometry} in either the parameter space or the activation space. For instance, classic weight-averaging (such as Model Soups \cite{wortsman22} or Task Arithmetic \cite{wang22, ilharco22}) and activation-blending (such as SABLE \cite{dauner23}) linearly interpolate parameters, activations, or projection matrices. While computationally cheap, this flat assumption is mathematically flawed when dealing with representation-space projection operators. 

In modern neural networks, representations are often projected onto low-dimensional subspaces to filter out task-irrelevant noise, align representations, or compress intermediate states \cite{bengio13, stoica23}. These low-rank task-specific projection bases or matrices represent $d$-dimensional subspaces embedded within a larger $D$-dimensional representation space $\mathbb{R}^D$. Mathematically, the set of all such $d$-dimensional subspaces is not a flat vector space; it is a curved, non-Euclidean manifold known as the \textbf{Grassmannian Manifold} $\mathcal{G}(d, D)$ \cite{absil08, edelman98}. 

When we linearly blend multiple orthogonal projection operators $\{P_k\}_{k=1}^K$ using flat ensembling weights $\{\alpha_k\}$ (where $\sum_{k=1}^K \alpha_k = 1$ and $\alpha_k \geq 0$) to form a merged projection $P_{\text{flat}} = \sum_{k} \alpha_k P_k$, the resulting operator violates the fundamental geometric properties of projection matrices. Specifically, $P_{\text{flat}}$ is almost never idempotent ($P_{\text{flat}}^2 \neq P_{\text{flat}}$). Its eigenvalues shrink strictly below $1$, causing what we formalize as \textbf{projected coordinate collapse} (see Section~\ref{sec:method}). As representations propagate sequentially through deep networks (e.g., the 14-layer Coordinate Sandbox analyzed in our experiments), this eigenvalue shrinkage compounds exponentially, decaying representation norms to zero and collapsing classification accuracy to a random baseline of $25.00\%$.

\begin{figure*}[t]
  \centering
  % We will copy the figure results/fig1.png or mock a placeholder path
  \includegraphics[width=0.85\textwidth]{../results/fig1.png}
  \caption{Joint Mean Classification Accuracy comparison under Orthogonal and Overlapping Manifolds across Homogeneous and Heterogeneous deployment streams. Traditional Uniform Merging experiences complete coordinate collapse ($25.00\%$) under overlapping manifolds. Our proposed \textbf{C-Lie-MM} resolves this collapse entirely, achieving stable, high-performance representation ensembling that is completely immune to heterogeneity collapse.}
  \label{fig:main_results}
\end{figure*}

To resolve this fundamental flaw, we introduce \textbf{Continuous Riemannian-Geometric Homotopical Model Merging via Grassmannian Geodesic Blending (C-Lie-MM)}, a mathematically rigorous framework designed specifically for curved subspace ensembling. Rather than interpolating projection matrices in a flat Euclidean space, C-Lie-MM treats each task-specific subspace as a point on the Grassmannian manifold $\mathcal{G}(d, D)$. We employ formal Riemannian-geometric operations to perform the ensembling. Specifically, we project the task-specific projection bases onto the tangent space of a reference subspace using the Grassmannian logarithm map $\log$. In this flat, vector-valued tangent space, we perform a weighted homotopic interpolation of the tangent matrices. Finally, we map the blended tangent vector back onto the manifold via the Grassmannian exponential map $\exp$. 

Because the exponential map maps the tangent space directly back onto the Grassmannian, the merged basis $Y_{\text{merged}}$ is guaranteed to be strictly orthogonal ($Y_{\text{merged}}^T Y_{\text{merged}} = I_d$). The resulting projection operator $P_{\text{merged}} = Y_{\text{merged}} Y_{\text{merged}}^T$ is thus guaranteed to be strictly symmetric, idempotent, and of rank $d$. This completely eliminates eigenvalue shrinkage and prevents projected coordinate collapse.

Our main contributions are summarized as follows:
\begin{enumerate}
    \item \textbf{Mathematical Formulation of Coordinate Collapse:} We mathematically formalize how flat linear blending of projection operators systematically violates the manifold geometry of the Grassmannian, proving that it leads to eigenvalue shrinkage and exponential representation decay in deep networks.
    \item \textbf{The C-Lie-MM Framework:} We propose Continuous Riemannian-Geometric Homotopical Model Merging via Grassmannian Geodesic Blending (C-Lie-MM), utilizing exact Riemannian log and exp mappings on $\mathcal{G}(d, D)$ to guarantee the geometric correctness of the merged projection operator.
    \item \textbf{Empirical Validation:} We evaluate C-Lie-MM inside a 14-layer Analytical Coordinate Sandbox simulation. C-Lie-MM completely resolves coordinate collapse, achieving $71.00\% \pm 2.37\%$ on orthogonal manifolds and $70.30\% \pm 4.01\%$ under severe manifold entanglement, outperforming flat ensembling baselines.
    \item \textbf{Immunity to Heterogeneity Collapse:} We demonstrate that C-Lie-MM is perfectly immune to heterogeneity collapse. By evaluating the Grassmannian geodesic barycenter sample-wise during the forward pass, C-Lie-MM maintains identical, high performance under both Homogeneous and Heterogeneous streaming workloads, presenting a major systems-level advantage.
    \item \textbf{Out-of-Distribution (OOD) Robustness:} We show that C-Lie-MM possesses an inherent geometric safeguard: under extreme uncertainty or out-of-distribution inputs, routing weights distribute evenly ($\alpha_k \to 1/K$), causing the ensembled tangent matrix to approach zero due to the properties of the Karcher mean centroid. The exponential map thus projects the representation directly onto the Karcher mean subspace $Y_0$, preserving shared semantic features and avoiding catastrophic collapse.
\end{enumerate}
